\chapter{Ingenier\'ia de requisitos}
\label{cap:requisitos}

Este cap\'itulo integra \texttt{docs/requisitos/SRS.md} (Especificaci\'on
de Requisitos de Software, con patr\'on ISO/IEC/IEEE~29148), las historias
de usuario, los casos de uso y \texttt{docs/requisitos/CHANGELOG-REQ.md}.
El SRS distingue expl\'icitamente requisitos \textbf{verificados},
\textbf{implementados} (sin prueba automatizada dedicada) y
\textbf{pendientes}; este cap\'itulo respeta esa misma distinci\'on: ning\'un
requisito que el SRS marca como pendiente se declara aqu\'i como cumplido.

\section{Priorizaci\'on MoSCoW}

Los requisitos se priorizan con la t\'ecnica MoSCoW (Must, Should, Could,
Won't). El n\'ucleo \emph{Must Have} corresponde al compromiso funcional
operativo de esta entrega: autenticaci\'on, autorizaci\'on por rol y por
propiedad, y CRUD de mascotas. Los m\'odulos de citas, consultas y
vacunas, incorporados durante la Unidad~IV, se formalizaron como
\texttt{REQ-F-023} a \texttt{REQ-F-025} (secci\'on~\ref{sec:req-unidad4}).

\section{Requisitos funcionales implementados y verificados}

Los requisitos \texttt{REQ-F-001} a \texttt{REQ-F-012} y \texttt{REQ-F-021}
cubren registro de usuario, login/logout/refresh por cookies,
autorizaci\'on por rol, perfil del usuario autenticado, y el CRUD completo
de mascotas con resumen por especie; todos con estado \textbf{verificado}
en el SRS, respaldados por pruebas automatizadas de controlador/servicio
(cap\'itulo~\ref{cap:pruebas-calidad}).

\subsection{Requisitos incorporados en la Unidad~IV (\texttt{REQ-F-023} a \texttt{REQ-F-025})}
\label{sec:req-unidad4}

Estos tres requisitos formalizan funcionalidades reales incorporadas al
backend durante la Unidad~IV, distintas de Citas y Consultas (que ya
ten\'ian requisito formal hist\'orico propio --- \texttt{REQ-F-015} y
\texttt{REQ-F-013} respectivamente --- y no se duplican aqu\'i):

\begin{itemize}[nosep]
  \item \textbf{REQ-F-023 --- Gesti\'on administrativa de usuarios}
    (\texttt{UsuarioController}/\texttt{UsuarioService}): un
    \texttt{ADMIN} lista, consulta, crea, actualiza y da de baja l\'ogica
    cuentas de cualquier rol, sin poder autoescalar el rol de su propia
    cuenta. Estado: \textbf{verificado} (\texttt{UsuarioControllerTest},
    16 pruebas).
  \item \textbf{REQ-F-024 --- Gesti\'on de vacunas} (\texttt{VacunaController}):
    registro, consulta (global, por mascota o por identificador),
    actualizaci\'on y baja l\'ogica de vacunas, con propiedad restringida
    para \texttt{ROLE\_DUENO}. Estado: \textbf{verificado}
    (\texttt{VacunaControllerTest}, 10 pruebas; colecci\'on Postman
    \texttt{BIOPET-Vacunas.postman\_collection.json}); la operaci\'on
    \texttt{GET /api/vacunas/mascota/\{mascotaId\}} no tiene prueba
    automatizada dedicada.
  \item \textbf{REQ-F-025 --- Consulta de informaci\'on externa de especies}
    (\texttt{ExternalApiController}): consulta de datos de una especie
    (nombre cient\'ifico, h\'abitat, dieta) desde un servicio externo (API
    Ninjas), con cach\'e Redis \emph{cache-aside} para evitar llamadas
    repetidas. Estado: \textbf{implementado} (sin prueba automatizada
    dedicada; no se marca ``verificado'' siguiendo la misma convenci\'on ya
    usada en el resto del SRS); evidencia emp\'irica manual de la
    comparaci\'on cach\'e fr\'ia/caliente en
    \rutalarga{docs/u4/evidencias/fred/redis-cache-comparacion.md}.
\end{itemize}

Rendimiento bajo carga de citas, consultas y vacunas no fue medido
espec\'ificamente con k6 (la medici\'on de rendimiento del
cap\'itulo~\ref{cap:pruebas-calidad}, secci\'on~\ref{sec:rendimiento-final},
se limita al endpoint de mascotas), limitaci\'on declarada en el
cap\'itulo~\ref{cap:amenazas-limitaciones-final}.

Las seis rutinas almacenadas de PostgreSQL (\texttt{fn\_resumen\_mascotas\_por\_especie},
\texttt{fn\_historial\_clinico\_mascota}, \texttt{fn\_reporte\_dashboard},
\texttt{fn\_siguiente\_numero\_ficha}, \texttt{sp\_actualizar\_estado\_citas\_masivas},
\texttt{sp\_registrar\_consulta\_validada}) respaldan estos m\'odulos con
acceso JPA formal (secci\'on~\ref{sec:acceso-jpa-formal}).

\section{Requisitos pendientes, heredados de la Entrega~1A}

Los requisitos \texttt{REQ-F-013} a \texttt{REQ-F-020} y \texttt{REQ-F-022}
permanecen, con honestidad, en estado \textbf{pendiente} o \textbf{parcial}
en el SRS, sin implementaci\'on de backend ni de frontend a la fecha de
este informe salvo lo explicitado:

\begin{table}[H]
  \centering
  \small
  \caption{Requisitos funcionales pendientes o parciales heredados de la Entrega 1A (fuente: \texttt{docs/requisitos/SRS.md}).}
  \label{tab:req-pendientes}
  \begin{tabular}{@{}p{2.0cm}p{6.4cm}p{1.6cm}p{2.0cm}@{}}
    \toprule
    Requisito & Descripci\'on & Prioridad & Estado \\
    \midrule
    REQ-F-013 & Registro/consulta de historial cl\'inico cronol\'ogico como m\'odulo independiente & Should & Pendiente \\
    REQ-F-014 & Registro de medicamentos prescritos & Could & Pendiente \\
    REQ-F-015 & Calendario interactivo de citas (m\'odulo b\'asico de citas ya existe; el calendario interactivo no) & Should & Pendiente \\
    REQ-F-016 & API de recepci\'on de telemetr\'ia IoT & Could & Pendiente \\
    REQ-F-017 & Recomendaciones cl\'inicas generadas por IA externa & Could & Pendiente \\
    REQ-F-018 & Facturaci\'on digital y registro de pagos & Should & Pendiente \\
    REQ-F-019 & Reportes estad\'isticos exportables (PDF/Excel) & Could & Pendiente \\
    REQ-F-020 & Auditor\'ia de operaciones del sistema & Should & Parcial \\
    REQ-F-022 & Notificaciones por correo electr\'onico & Could & Pendiente \\
    \bottomrule
  \end{tabular}
\end{table}

\textbf{Nota de honestidad sobre \texttt{REQ-F-015} y \texttt{REQ-F-013}.}
El backend expone \texttt{CitaController} y \texttt{ConsultaController}
reales, con endpoints funcionales y respaldados por las rutinas
almacenadas de la Unidad~IV (secci\'on~\ref{sec:req-unidad4}); sin embargo,
el SRS mantiene \texttt{REQ-F-013} (historial cl\'inico cronol\'ogico
completo, con consulta longitudinal) y la parte de \emph{calendario
interactivo} de \texttt{REQ-F-015} como pendientes, porque esas
funcionalidades espec\'ificas --- m\'as all\'a del CRUD b\'asico de
consultas y citas ya implementado --- no est\'an completas ni tienen
frontend dedicado. Este informe reporta ambos estados tal como los declara
el SRS vigente, sin resolver la aparente tensi\'on ampliando el alcance
por cuenta propia: se trata como lo que es, una actualizaci\'on parcial del
SRS que a\'un no reclasific\'o expl\'icitamente cada subrequisito heredado
frente al c\'odigo nuevo de la Unidad~IV, y se documenta como limitaci\'on
de consistencia documental en el cap\'itulo~\ref{cap:amenazas-limitaciones-final}.

\texttt{REQ-F-020} (auditor\'ia) se marca \textbf{parcial}, no
\textbf{pendiente}: el registro de eventos de autenticaci\'on
(\texttt{AUTH\_AUDIT}, cap\'itulo~\ref{cap:seguridad}) s\'i existe y est\'a
verificado, pero no cubre auditor\'ia gen\'erica de todas las operaciones
CRUD sobre mascotas ni de los m\'odulos de citas/consultas/vacunas.

\section{Requisitos no funcionales relevantes}

\begin{table}[H]
  \centering
  \small
  \caption{Requisitos no funcionales relevantes citados en este informe (fuente: \texttt{docs/requisitos/SRS.md}, secci\'on 3.2).}
  \label{tab:req-nf-relevantes}
  \begin{tabular}{@{}p{2.0cm}p{6.6cm}p{2.6cm}@{}}
    \toprule
    Requisito & Descripci\'on & Estado \\
    \midrule
    REQ-NF-001 & Rendimiento del listado de mascotas ($\leq$200~ms p95 caliente, $\leq$500~ms fr\'ia) & Verificado \\
    REQ-NF-002 & Comunicaci\'on cifrada obligatoria (HTTPS/TLS) & Verificado (TLS de desarrollo) \\
    REQ-NF-003 & Expiraci\'on configurable de tokens JWT & Verificado \\
    REQ-NF-004 & Claims est\'andar del JWT (RFC~7519) & Verificado \\
    REQ-NF-005 & Interfaz responsiva (320--1440px) / Lighthouse & Verificado (secci\'on~\ref{sec:lighthouse-final}) \\
    REQ-NF-006 & Compatibilidad de navegadores (Chrome/Firefox/Edge) & Pendiente de evidencia archivada \\
    REQ-NF-008 & Cifrado de contrase\~nas (BCrypt) & Verificado \\
    REQ-NF-009 & Registro de eventos de autenticaci\'on (OWASP A09) & Verificado \\
    REQ-NF-010 & Limitaci\'on de intentos de login (OWASP A07) & Verificado \\
    REQ-NF-011 & Arquitectura en capas (\texttt{controller}/\texttt{service}/\texttt{repository}) & Verificado \\
    REQ-NF-012 & Cach\'e del listado de mascotas con TTL externo & Verificado parcialmente (TTL confirmado; sin hit ratio medido) \\
    REQ-NF-013 & Estrategia h\'ibrida de acceso a datos (JPA + procedimientos almacenados) & Verificado (cap\'itulo~\ref{cap:arquitectura}) \\
    \bottomrule
  \end{tabular}
\end{table}

El requisito de interfaz responsiva/calidad web (\texttt{REQ-NF-005}) est\'a
\textbf{verificado}: las cuatro categor\'ias de Lighthouse (Performance,
Accessibility, Best Practices, SEO) cumplen su umbral en la corrida del
2026-08-18, en ambos perfiles mobile y desktop
(secci\'on~\ref{sec:lighthouse-final}) --- superando el estado
``verificado parcialmente'' que declaraba una revisi\'on anterior del SRS,
fechada antes de esa corrida. \texttt{REQ-NF-012} (hit ratio de cach\'e
Redis) permanece verificado solo parcialmente: existe evidencia de TTL y
de existencia de clave bajo carga, pero no una medici\'on expl\'icita de
\texttt{keyspace\_hits}/\texttt{keyspace\_misses} a lo largo del tiempo.
No existe, en el SRS actual, un requisito no funcional formal dedicado
espec\'ificamente a la usabilidad medida con SUS: la medici\'on SUS
(secci\'on~\ref{sec:sus-final}) responde a un criterio de evaluaci\'on de
la Gu\'ia de la Entrega Final, no a un \texttt{REQ-NF} numerado del SRS.

\section{Trazabilidad y estado de la validaci\'on autom\'atica}

\texttt{scripts/validate-traceability.sh} valida de forma cruzada que todo
requisito del SRS tenga fila en \texttt{docs/trazabilidad/matriz.csv} con
historia de usuario y caso de uso asociados. El detalle completo de la
matriz se presenta en el cap\'itulo~\ref{cap:trazabilidad-final}.

\section{Marco INCOSE}

La redacci\'on de cada requisito sigue el patr\'on ``El sistema
deber\'a\ldots'' con entradas, resultado esperado y criterios de
aceptaci\'on verificables, alineado con las caracter\'isticas de calidad de
requisitos de la \emph{INCOSE Guide to Writing Requirements}~\cite{incose2023requirements}
(\texttt{docs/checklists/incose2023-req.md}): en particular
\emph{Unambiguous} y \emph{Verifiable}, aplicadas expl\'icitamente al
cerrar la ambig\"uedad hist\'orica de \texttt{REQ-F-017} (observaci\'on
OBS-04, ya cerrada).

\section{Cambios entre versiones}

\texttt{docs/requisitos/CHANGELOG-REQ.md} documenta la evoluci\'on del SRS
entre la Entrega~1A, la Entrega~1B, la Tercera Entrega y la Entrega Final,
incluyendo la incorporaci\'on de \texttt{REQ-F-022} (recuperaci\'on formal
de RF-07, cierre de OBS-02) y de \texttt{REQ-F-023} a \texttt{REQ-F-025}
(Unidad~IV). Ning\'un identificador de requisito se reutiliz\'o para dos
funcionalidades distintas.
